\textbf{Project:} \emph{From Cognitive Architectures to Production
Agentic Systems: Understanding CoALA, Fan \& Lan, and Engineering
Complex Agents with Java}\\
\textbf{Status:} Authoritative working specification. Read this file
before outlining, drafting, revising, or packaging any chapter.

\subsection{1. Approval and delivery
workflow}\label{1-approval-and-delivery-workflow}

\begin{enumerate}
\def\labelenumi{\arabic{enumi}.}
\tightlist
\item
  Every chapter receives its own detailed outline and its own approval
  before chapter prose is written.
\item
  Work in batches of two chapters for drafting efficiency, but retain
  chapter-by-chapter outline approval. One chapter\textquotesingle s
  approval does not approve the next chapter\textquotesingle s outline.
\item
  For each batch: prepare the two outlines as individual review PDFs;
  wait for approval of each; draft only the approved chapters; perform a
  source-fidelity and cross-chapter consistency review; build and
  visually inspect individual chapter PDFs and the combined book PDF.
\item
  Do not silently treat a proposed outline, a prior
  chapter\textquotesingle s approval, or a request to continue as
  approval of an unapproved outline.
\item
  Deliver chapter content as individual PDFs. Do not deliver Markdown
  chapters, ZIP archives, or combined PDFs unless the author explicitly
  requests those formats. Repository source files remain in the project
  repository.
\item
  Keep a durable status entry in \texttt{README.md} and preserve
  approved outlines, rule versions, and delivery checkpoints in
  \texttt{review/}.
\end{enumerate}

\subsection{2. Source discipline and theoretical
scope}\label{2-source-discipline-and-theoretical-scope}

The book has two formal theoretical foundations only:

\begin{itemize}
\tightlist
\item
  \textbf{{[}CoALA{]}} Sumers, Yao, Narasimhan, and Griffiths,
  \emph{Cognitive Architectures for Language Agents}, Transactions on
  Machine Learning Research (2024), with the original arXiv manuscript
  as the earlier version.
\item
  \textbf{{[}Fan \& Lan{]}} Fan and Lan, \emph{From Cognitive
  Architectures to Language Agents: A Mechanism-Level Review of Lineage,
  Convergence, and Migration Gaps} (2026).
\end{itemize}

Use only the two formal source labels \texttt{{[}CoALA{]}} and
\texttt{{[}Fan\ \&\ Lan{]}}. Use the shared
\texttt{\textbackslash{}CoALACite\{...\}} and
\texttt{\textbackslash{}FanLanCite\{...\}} footnote macros. Do not
invent additional formal theory labels.

Part I explains CoALA on its own terms. Never use Fan \& Lan to
reinterpret, extend, or retroactively change what CoALA says. Part II
treats Fan \& Lan as a separate paper and distinguishes its explicit
claims, hypotheses, analysis, and boundaries. Part III\textquotesingle s
Java systems are engineering applications, not claims that either paper
prescribes those implementations.

Do not introduce Tree of Thoughts, RAP, or any other research framework
as an additional theoretical foundation. Other works may be named only
when they are examples explicitly discussed or cited within the primary
paper being explained; do not turn those works into parallel theory
chapters or claim their mechanisms as CoALA/Fan \& Lan mechanisms.

For every substantive theoretical claim, identify the exact supporting
paper section, figure, table, or appendix during research and preserve
that trace in the prose and footnotes. If a claim is an explanation,
interpretation, illustrative analogy, or engineering recommendation
rather than an explicit paper statement, mark and phrase it accordingly.
Do not make unsupported claims sound like quotations or direct paper
statements.

\subsection{3. Teaching style and chapter
construction}\label{3-teaching-style-and-chapter-construction}

\begin{itemize}
\tightlist
\item
  Write for a technically curious beginner with basic programming
  knowledge; do not assume prior knowledge of cognitive architectures or
  agent research.
\item
  Use the sequence \textbf{problem first → concrete intuition → term →
  precise mechanism → example → boundary}. Do not dump a glossary of
  definitions before the reader needs the terms.
\item
  Explain every new term at first meaningful use. Distinguish nearby
  concepts explicitly.
\item
  Use complete, connected explanations, not compressed notes or lists of
  unexplained labels. Tables and bullets support explanations; they do
  not replace them.
\item
  Develop concepts progressively. Reuse the Product X payment-service
  incident as the book-created running example where useful, but
  identify it as a \textbf{{[}Book example{]}} and never imply it comes
  from either paper.
\item
  Make examples technically plausible and internally consistent. Do not
  add invented product details and then use them as factual evidence.
\item
  Each chapter should have a clear chapter-at-a-glance orientation,
  learning outcomes, a logical section progression, selective worked
  examples, source/boundary notes, a recap, and review questions or
  exercises when pedagogically useful.
\item
  Avoid repeating material already taught. Preview later concepts
  briefly and point to their chapter; do not teach later chapters
  prematurely.
\item
  Keep chapters substantial enough to explain the subject properly.
  There is no fixed page-count ceiling; completeness, clarity, accuracy,
  and teaching value govern length.
\end{itemize}

\subsection{4. Part I scope boundaries}\label{4-part-i-scope-boundaries}

\begin{itemize}
\tightlist
\item
  Chapter 1 motivates why a cognitive architecture is useful and
  introduces the CoALA map at a high level.
\item
  Chapter 2 explains the architecture as a whole and the relationships
  among memory, action space, and the decision procedure.
\item
  Chapter 3 explains the memory model in depth: working memory and
  long-term episodic, semantic, and procedural memory.
\item
  Chapter 4 explains internal actions: retrieval, reasoning, and
  learning. Keep external grounding as a boundary and defer its full
  explanation to Chapter 5.
\item
  Chapter 5 explains external actions and grounding.
\item
  Chapters 6--8 develop planning, proposal/evaluation/selection, and
  interleaving/iteration. Do not duplicate the full decision procedure
  in Chapters 3 or 4.
\end{itemize}

\subsection{5. Shared LaTeX system and
preamble}\label{5-shared-latex-system-and-preamble}

\begin{itemize}
\tightlist
\item
  \texttt{Tex/preamble.tex} is the one shared authorial preamble. It is
  loaded once by \texttt{Tex/main.tex}; chapter fragments must never
  load it.
\item
  The latest author-supplied preamble supersedes earlier preferences for
  trim size and fonts: its current specification is a 7 × 10 inch,
  two-sided, open-right book using TeX Gyre Termes, TeX Gyre Heros, and
  TeX Gyre Cursor. The author-supplied preamble is the design baseline.
  It may be changed only when a demonstrable improvement is needed for
  consistent typography, accessibility, rendering, pagination, or figure
  placement. Preserve existing design intent and explain nontrivial
  changes in the review notes.
\item
  Use the \texttt{book} class, XeLaTeX, \texttt{fontspec}, the
  author\textquotesingle s selected fonts and page layout, shared
  colors, heading styles, headers/footers, \texttt{listings},
  \texttt{graphicx}, captions, and reusable teaching environments from
  the preamble.
\item
  Use \texttt{Tex/shared-macros.tex} for book-wide commands, citations,
  glossary, index, and reusable macros/environments. Do not create
  chapter-local versions of shared commands or boxes.
\item
  Shared environments include \texttt{learningbox}, \texttt{tipbox},
  \texttt{warningbox}, \texttt{notebox}, \texttt{takeawaybox},
  \texttt{chapterrecap}, and \texttt{reviewquestions}. Use them
  consistently and selectively. Chapter 1 and Chapter 2 must follow the
  same semantic box policy. Boxes are for learning objectives, genuinely
  useful tips, important cautions, clarifications, takeaways, recaps,
  and questions---not decoration or wrapping ordinary paragraphs.
\item
  Keep code examples in \texttt{listings}; use the shared language
  definitions and formatting. Avoid raw code screenshots.
\item
  Use semantic macros and labels consistently. Every figure has a
  caption and stable label; every cross-reference resolves.
\end{itemize}

\subsection{6. Diagrams and figures: Mermaid is
authoritative}\label{6-diagrams-and-figures-mermaid-is-authoritative}

All conceptual diagrams are authored as Mermaid text and rendered to PNG
by CI before LaTeX compilation.

\begin{itemize}
\tightlist
\item
  Store editable \texttt{.mmd} sources under
  \texttt{Tex/diagrams/source/partN/chapterNN/}.
\item
  Store shared Mermaid theme and configuration in
  \texttt{Tex/diagrams/mermaid.config.json} (and related shared renderer
  configuration only when needed).
\item
  Pin Mermaid CLI in \texttt{tools/mermaid/package.json} and its
  lockfile. Never use an unpinned
  \texttt{npx\ @mermaid-js/mermaid-cli@latest} in CI.
\item
  CI uses Node.js 22.13+ (the repository currently configures Node 24),
  runs \texttt{npm\ ci}, and renders every \texttt{.mmd} source into a
  high-resolution PNG under the matching path in
  \texttt{Tex/diagrams/generated/} before XeLaTeX starts.
\item
  The shared render script must be deterministic and recursive, handle
  spaces safely, fail on invalid Mermaid syntax, and generate a PNG for
  every source. CI validates PNG integrity, dimensions/resolution,
  matching LaTeX references, and missing/orphan generated files.
\item
  Use one shared Mermaid visual theme derived from the book colors and
  typography. Use \texttt{accTitle} and \texttt{accDescr} in each
  diagram source; also provide meaningful captions and nearby
  explanatory prose because the book embeds PNGs.
\item
  Use Mermaid for flowcharts, architecture maps, decision cycles,
  timelines, state diagrams, and other conceptual illustrations. Do not
  use TikZ, hand-positioned LaTeX nodes, ASCII-art diagrams, or custom
  LaTeX diagram libraries for book diagrams.
\item
  LaTeX only places the generated image with \texttt{graphicx} through
  the shared \texttt{\textbackslash{}BookDiagram} macro. The macro owns
  consistent width/height constraints, caption and label placement;
  Mermaid owns node arrangement, routing, spacing, and diagram layout.
\item
  Keep diagram labels concise. If a diagram becomes too dense, split it
  into multiple focused diagrams rather than shrinking text until it is
  unreadable.
\item
  Verify the actual CI-generated PNGs and final PDF at printed size. A
  mocked renderer or placeholder PNG can test the LaTeX integration
  layer only; it must never be reported as proof that real Mermaid
  output rendered correctly.
\item
  Generated PNGs are build artifacts (not source of truth). CI uploads
  the generated PNG directory as a separate review artifact.
\end{itemize}

\subsection{7. References, glossary, index, and
terminology}\label{7-references-glossary-index-and-terminology}

\begin{itemize}
\tightlist
\item
  Bibliography is managed by \texttt{biblatex} and Biber, with
  footnote-style citations as selected by the author.
\item
  Use \texttt{\textbackslash{}CoALACite\{...\}} and
  \texttt{\textbackslash{}FanLanCite\{...\}} for paper citations; do not
  hard-code one-off citation formats in chapter files.
\item
  Add glossary entries to \texttt{Tex/glossary/terms.tex} and index
  entries consistently. Avoid duplicating terms or changing a
  term\textquotesingle s meaning between chapters.
\item
  Keep terminology stable across chapters. When a concept is revisited,
  deepen it rather than redefining it differently.
\end{itemize}

\subsection{8. Quality assurance and release
gates}\label{8-quality-assurance-and-release-gates}

Before delivering a chapter or batch:

\begin{enumerate}
\def\labelenumi{\arabic{enumi}.}
\tightlist
\item
  Check all claims against the relevant primary-paper sections and
  figures.
\item
  Check no claims from one paper are attributed to the other.
\item
  Check the outline was explicitly approved.
\item
  Check chapter scope and avoid duplication with preceding/following
  chapters.
\item
  Compile with XeLaTeX through \texttt{latexmk}, using the
  repository\textquotesingle s GitHub Actions TeX Live workflow.
\item
  Run the repository\textquotesingle s scaffold, Mermaid, and PDF
  quality checks.
\item
  Inspect the PDF visually after rendering pages to images. Check
  diagrams at their printed size, box consistency, page breaks, headers,
  captions, references, footnotes, and overfull boxes.
\item
  Extract PDF text and verify chapter titles, headings, labels,
  metadata, and searchability.
\item
  Report which validation ran locally and which requires a real GitHub
  Actions execution. Never describe a mocked or substitute-font test as
  a full production build.
\item
  Deliver chapter PDFs individually. Include a build note for any known
  rendering limitation; do not claim accessibility conformance beyond
  the selected standard PDF target.
\end{enumerate}

\subsection{9. Current project state
(2026-11-14)}\label{9-current-project-state-2026-11-14}

\begin{itemize}
\tightlist
\item
  Chapters 1 and 2 have been drafted and integrated.
\item
  Mermaid-based diagram generation was introduced for Chapters 1 and 2;
  the GitHub Actions run is the authoritative real-render verification
  point.
\item
  The next outline batch is Chapters 3 and 4. Both outlines must be
  reviewed and approved before either chapter\textquotesingle s prose is
  written.
\item
  Chapters 5--30 remain placeholders until their own outlines are
  approved.
\end{itemize}
